\clearpage
\section{Quản lý phiên bản và tính nhất quán dữ liệu}
\label{sec:dataset-registry}

Bộ dữ liệu được lưu trong cơ sở dữ liệu SQLite với các phiên bản bất biến. Mục đích là bảo đảm mỗi kết quả thực nghiệm truy ngược được đúng quỹ đạo, nhãn và cửa sổ quan sát đã sử dụng. SQLite phù hợp kho nghiên cứu cục bộ, không cần dịch vụ máy chủ; các thao tác nhập được thực hiện trong một giao dịch và bật kiểm tra khóa ngoại trên từng kết nối.~\citep{sqlitetransactions,sqliteforeignkeys,sqliteuses} Đây là thành phần quản lý thực nghiệm, không phải một cơ chế bảo vệ vị trí mới.

\subsection{Tổ chức dữ liệu và lịch sử thay đổi}

Một phiên bản chứa các nhóm mô phỏng, mỗi nhóm thuộc đúng một phần dữ liệu học, chọn cấu hình hoặc kiểm thử/xác nhận. Mỗi phiên di chuyển thuộc một nhóm và giữ riêng mã người, thiết bị, xe vật lý. Chuỗi FCD được lưu theo khóa gồm phiên bản, mã phiên và chỉ số mẫu; thời điểm phải tăng nghiêm ngặt. Vĩ độ, kinh độ, vận tốc, làn, cạnh đường, vị trí dọc làn và góc hướng được giữ nguyên từ bản xuất SUMO.

Danh mục kịch bản được lưu độc lập với các mẫu đã sinh, nên ca chưa có dữ liệu không biến mất khỏi mẫu số. Một bản ghi kịch bản tham chiếu tới một hoặc nhiều phiên và danh sách chỉ số quan sát được phép. Khóa ngoại bảo đảm chỉ số trỏ tới mẫu có thật, phiên thuộc đúng nhóm và bản ghi không bị gán sang phần dữ liệu khác. Nhãn đặc thù, điều kiện chấp nhận, cấu hình sinh và nguồn gốc được lưu kèm; tính đúng đắn vật lý của các trường này vẫn cần bộ kiểm chứng kịch bản.

Khi nhập một bản mới, hệ thống kiểm tra cấu trúc, thời gian, quan hệ định danh, chỉ số và số lượng khai báo; sau đó đối chiếu bản xuất lại với nội dung gốc. Chỉ khi toàn bộ bước thành công mới khóa phiên bản và ghi vào bảng \texttt{update\_log}. Mỗi mục lịch sử gồm phiên bản trước/sau, người thực hiện, thời điểm nhập, lý do, mã băm và số dòng ở các lớp dữ liệu. Thời điểm nhật ký là thời gian quản lý dữ liệu, không phải đồng hồ mô phỏng của chuyến đi.

Thay đổi quỹ đạo, nhãn hoặc cách chia dữ liệu phải tạo phiên bản mới, không sửa trực tiếp bản đã khóa. Ghi dữ liệu và nhật ký cùng một giao dịch giúp tránh tình trạng nhập dở hoặc dữ liệu đã đổi nhưng không có lịch sử. Hai tiến trình cùng cập nhật phải khai báo phiên bản cha dự kiến; tiến trình có trạng thái cũ bị từ chối nếu tiến trình khác đã công bố bản mới. Đây là chuỗi phiên bản tuyến tính cho nghiên cứu cục bộ, chưa phải hệ thống cộng tác nhiều nhánh.

\subsection{Cố định đầu vào và ranh giới truy cập}

Mỗi lần chạy cần cố định cặp $(v,h_v)$, trong đó $v$ là mã phiên bản và $h_v$ là SHA-256 của nội dung JSON chuẩn hóa: sắp thứ tự khóa, giữ thứ tự danh sách và giá trị, không phụ thuộc thụt lề. Mã băm theo byte của tệp gốc được lưu riêng; hai mã không được dùng thay nhau. Không dùng tên \emph{mới nhất} như đầu vào ngầm định, vì tên đó có thể trỏ sang dữ liệu khác sau một lần cập nhật.

Đối với benchmark đã công bố, giữ nguyên bản JSON và mã băm ban đầu để phát lại đúng giao thức cũ. Kho SQLite cung cấp phiên bản đã khóa và luồng quan sát cho các bước phát triển tiếp theo; việc chuyển nơi lưu không tạo một bộ dữ liệu khoa học mới. Kiểm tra di chuyển v1 và v2 đối chiếu 116.152 điểm, 303 bản ghi và 468 luồng phiên được phép. Tất cả 8.466 lượt quan sát khớp bản gốc. Một điểm có thể được dùng bởi nhiều kịch bản, nên số lượt quan sát không phải số điểm hay số người độc lập.

Cơ sở dữ liệu là kho đáp án của bộ đánh giá, \textbf{không phải dữ liệu gửi tới LSP}. Giao diện đầu vào thiết bị chỉ xuất thời gian tương đối, tọa độ và loại truy vấn trong cửa sổ được phép. S8 giữ chênh lệch thời gian giữa hai phiên bằng cùng mốc quan sát; S6.C phân biệt sáu phiên lịch sử với phiên hiện tại. Các quan sát thật này vẫn phải qua cơ chế bảo vệ trước khi đối thủ được xem. Tương lai, định danh nội bộ, làn đường và nhãn đánh giá không được ghép vào bản công bố.

\subsection{Điều kiện dùng một kịch bản để benchmark}

Cần phân biệt ba mức bằng chứng: \emph{có mẫu hợp lệ}, \emph{đủ dữ liệu cho giao thức đánh giá}, và \emph{đã đo hiệu quả bảo vệ}. Cơ sở dữ liệu hỗ trợ mức đầu về tính nhất quán và đếm mẫu, nhưng không tự suy ra hai mức sau. Để kiểm tra độ bao phủ, lấy tích của danh mục 30 ca với ba phần dữ liệu v2 rồi đếm cả các tổ hợp bằng không. Một ca chỉ đủ điều kiện cho giao thức có học/chọn đối thủ khi có dữ liệu phù hợp trong các phần cần thiết, nhãn đối chứng đủ và số nhóm độc lập được công bố.

Chẳng hạn S5.C có số bản ghi học/chọn cấu hình/xác nhận lần lượt là $0/0/1$. Việc ca này xuất hiện trong hợp bộ dữ liệu không cho phép tuyên bố đã benchmark được đối thủ cần huấn luyện. Đối thủ định sẵn có thể chạy như phép thử chẩn đoán trên mẫu đơn, nhưng không hỗ trợ kết luận tổng quát. Khi bổ sung dữ liệu, cần sinh thêm nhóm theo điều kiện kịch bản đã định trước, không chọn mẫu bằng điểm số bảo vệ thuận lợi. Tập xác nhận đã được dùng để sửa thuật toán phải được xem là dữ liệu phát triển cho vòng sau; cần tập xác nhận mới.

Khóa ngoại, giao dịch và mã băm không thay thế kiểm tra tuyến SUMO, điều kiện tấn công, tính đại diện của giao thông hoặc rò rỉ thông tin qua giao thức. Nhật ký và các ràng buộc chống sửa nhầm trong quy trình thông thường, không chống quản trị viên có toàn quyền thay cơ sở dữ liệu và mã kiểm tra. Dữ liệu liên quan trong cùng nhóm không trở thành các quan sát độc lập chỉ vì được lưu thành nhiều dòng. Các giới hạn này phải được giữ khi diễn giải kết quả và mở rộng sang nhiều thành phố.

\subsection{Phiên bản xác nhận cho vòng tối ưu phủ dịch vụ}

Phiên bản v3 giữ nguyên nội dung bốn nhóm phát triển 101--104 của v2 và bổ sung bốn nhóm SUMO mới 201--204. Hai nhóm 105/106 đã được xem kết quả không tham gia xác nhận vòng mới. Mã định dạng bản ghi vẫn là v2 vì cấu trúc không đổi; mã phiên bản dữ liệu trong kho là \texttt{urban-scenarios-v3}. Tám nhóm chứa 176 phiên di chuyển, 124.852 mẫu FCD và 269 bản ghi kịch bản. Hợp dữ liệu có đủ 30 ca, nhưng từng nhóm chỉ đạt 24--30 ca; không diễn giải độ phủ hợp thành đủ dữ liệu cho mọi phép thử.

Các điều kiện sinh và loại mẫu của v2 được giữ nguyên. Ở S1--S3, mỗi ca có hai nhóm học, hai nhóm chọn cấu hình và bốn nhóm xác nhận. Các thiếu hụt ở phần khác vẫn hiện trong danh mục: S5.C chỉ có một bản ghi và không có mẫu học/chọn cấu hình. Bản ghi thiếu điều kiện không được sửa tọa độ hoặc đổi nhãn để làm đẹp độ phủ.

Bộ kiểm chứng đối chiếu lại 72 lượt SUMO theo ngày, gồm cả nguồn của bốn nhóm được giữ lại; toàn bộ 124.852 mẫu khớp FCD gốc. Kiểm tra tuyến gồm 9.080 lần chuyển cạnh ngoài. Các số này kiểm chứng nguồn gốc và điều kiện mô phỏng, chưa chứng minh tính đại diện cho dân cư đô thị.

Vòng thí nghiệm mới đọc trực tiếp kho SQLite với mã phiên bản và mã băm nội dung cố định. Prior chỉ lấy các phiên học; đầu vào mỗi cơ chế chỉ gồm cửa sổ quan sát được phép. Các thí nghiệm trước vẫn dùng ảnh chụp và mã băm cũ để phát lại. Việc thêm v3 tạo một dòng nhật ký mới và không thay nội dung hai phiên bản đã khóa.

\subsection{Tuyến phụ trợ cho học đối thủ suy luận}
\label{sec:auxiliary-shadow-data}

Học đối thủ từ nhiều đầu ra ngẫu nhiên trên cùng vài vị trí chưa tạo ra độ phủ không gian rộng. Vì vậy, kho bổ sung phiên bản \texttt{urban-shadow-v1}, độc lập về mục đích với bộ ca S1--S10. Phiên bản này chứa 80 nhóm tuyến SUMO, 160 phiên, 175.211 mẫu FCD nguyên gốc và 320 cửa sổ quan sát. Có 64 nhóm dùng học đối thủ, 16 nhóm giữ lại để kiểm tra chuyển giao; ba phiên bản dữ liệu cũ không thay đổi.

Tuyến được lập từ bản đồ đường công khai, không dựa vào vị trí thật của nhóm cần đánh giá. Chọn ô xuất phát 120 m không hoàn lại; đường đi cơ sở dài 2,5--4 km, có ít nhất 12 cạnh ngoài nút giao. Mỗi nhóm gồm một chuyến di chuyển và một chuyến quay lại cùng điểm dừng, với hai lần dừng 180 giây. Điểm dừng giữa các nhóm phải thuộc làn và ô 1 m khác nhau. Đây là lấy mẫu có kiểm soát độ phủ, không phải mô phỏng phân bố cư dân hoặc nhu cầu giao thông đã hiệu chuẩn. SUMO thực thi chuyển động ở bước 1 giây, vận tốc tối đa 8 m/s, không dịch chuyển tức thời và dừng trên làn.~\citep{sumotrip}

\begin{table}[htbp]
\centering\small
\caption{Bốn dạng cửa sổ phụ trợ; mỗi dạng có 64 nhóm học và 16 nhóm giữ lại. Đây là lịch quan sát, không phải bốn kịch bản đe dọa mới.}
\label{tab:auxiliary-profiles}
\begin{tabularx}{\textwidth}{p{2.6cm}X}
\toprule
Dạng quan sát & Quy tắc lấy từ FCD đã thực hiện\\
\midrule
Di chuyển, 20 s & Tối đa 12 điểm, cách nhau 20 giây trên chuyến cơ sở.\\
Di chuyển, 60 s & Tối đa 12 điểm, cách nhau 60 giây từ đầu chuyến cơ sở.\\
Dừng, 5 s & 12 điểm cách nhau 5 giây trong lần dừng thứ nhất.\\
Quay lại, 5 s & Sáu điểm ở mỗi lần dừng; giữa hai lần quan sát có khoảng trống trên 120 giây và chuyển động thực.\\
\bottomrule
\end{tabularx}
\end{table}

Các cửa sổ học có 2.873 quan sát, 1.345 tọa độ khác nhau trên 587 ô 120 m; tập giữ lại có 714 quan sát, 322 tọa độ khác nhau trên 208 ô. Số điểm khác nhau không phải số người độc lập: nhiều cửa sổ trong một nhóm vẫn liên quan. Không có tuyến hoàn chỉnh hoặc cửa sổ trùng chính xác với v3, cũng không có cửa sổ trùng giữa các nhóm phụ trợ. Tuy nhiên, các tuyến mới dùng chung 901 cạnh với v3; khoảng cách từ nhãn giữ lại đến nhãn học gần nhất có trung vị 30,47 m và phân vị 95 là 159,59 m. Đây là kiểm tra trên cùng bản đồ, không phải tổng quát hóa sang vùng địa lý mới.

Kiểm tra độc lập đối chiếu toàn bộ FCD với XML gốc, tuyến tới đích, thời gian dừng, liên tục thời gian và vị trí trên làn trước khi ghi kho. Hai thiết kế ban đầu bị loại do chế độ đỗ làm lệch điểm khỏi làn và cửa sổ dừng trùng giữa các nhóm; quy tắc được sửa trước khi học hoặc chấm đối thủ. Dữ liệu không được thêm nhiễu hay chỉnh tọa độ để vượt kiểm tra. Các cửa sổ phụ trợ dùng cho phép kiểm tra ở Mục~\ref{sec:expanded-shadow-experiment}; chúng không thay thế các điều kiện hợp lệ riêng của bộ ca chính.
